% !TEX TS-program = pdflatexmk
% !TEX program = pdflatex
% Typeset in TeXShop, or run make problem HW=hw05 from the repository root.
% Based on problems/template.tex; the template is unchanged.
% Source: MA-GY-6213-hw5.pdf, pages 1--2.
% MA-GY-6213, Homework No. 5, due Friday 10/16/2026.
% Page 1 contains Q1--Q6 without numbered subparts.
% Page 2 contains Q7--Q10; Q8--Q10 each have parts (a)--(b).
% Ten questions; the assignment supplies no solutions.
% The Week 5 summary is internal coverage metadata, not another exercise.
\documentclass[nohyper,nobib,nofonts]{tufte-handout}
% Resolve shared inputs from either the repository root or problems/.
\newcommand{\analysisroot}{..}
\IfFileExists{preamble.tex}{\renewcommand{\analysisroot}{.}}{}
\makeatletter
\edef\input@path{{\analysisroot/}}
\makeatother
% Shared packages and typography for the book and homework handouts.
% Handouts set this directory when TeXShop compiles from problems/.
\providecommand{\analysisroot}{.}
\usepackage[T1]{fontenc}
% Give titletoc a separate part number before nameref wraps \@part.
\usepackage{etoolbox}
\makeatletter
\@ifclassloaded{tufte-book}{%
  \patchcmd{\@part}
    {\thepart\hspace{1em}}
    {\protect\numberline{\thepart}}
    {}
    {\PackageError{analysis-notes}{Could not format part contents labels}
      {Check the book class definition of \string\@part.}}
}{}
\makeatother

\usepackage{nameref}
\usepackage{amsmath,amssymb,amsthm}
\usepackage{aliascnt}
\usepackage{graphicx}
\usepackage{tikz}
\usetikzlibrary{arrows.meta}
\usepackage{booktabs}
\usepackage{multicol} % Formula sheet and multicolumn specimen examples.
\usepackage{needspace} % Keep a section opening together when space is short.
% Packages used by the temporary Tufte specimen content.
\usepackage{fancyvrb,xspace,units,makeidx}
\fvset{fontsize=\small}
\usepackage{url}
\usepackage{lipsum} % Remove after all placeholder prose has been replaced.
% AMS numeric references, alphabetised by author, using the AMS's own style.
% Keep the resource path relative so handouts also build from problems/.
\usepackage[numbers,sort&compress]{natbib}
\bibliographystyle{amsplain}
\newcommand{\analysisbibliography}{%
  % A leading ./ prevents BibTeX from using latexmk's adjusted search path
  % when it changes into build/. Direct problems/ builds need ../ instead.
  \ifthenelse{\equal{\analysisroot}{.}}%
    {\bibliography{bibliography}}%
    {\bibliography{../bibliography}}%
}

\graphicspath{{\analysisroot/figures/}}
\setkeys{Gin}{width=\linewidth,totalheight=\textheight,keepaspectratio}
\setcounter{secnumdepth}{1}

% Retain the borrowed title/contents design only for the book.
% Adapted from Kevin Godby's Tufte-LaTeX title/contents example:
% https://groups.google.com/forum/#!topic/tufte-latex/ujdzrktC1BQ
\makeatletter
\@ifclassloaded{tufte-book}{%
  \makeindex
  % Leave room for two-digit appendix figure numbers such as C.13.
  \patchcmd{\@tufte@lof@line}{\@tempdima 2.0em}{\@tempdima 3.0em}
    {}
    {\PackageError{analysis-notes}{Could not widen list-of-figures numbers}
      {Check the Tufte list-of-figures definition.}}
  \setcounter{tocdepth}{0} % Parts and chapters only.
  % Title-only chapter headings; retain counters for sections and references.
  \titleformat{\chapter}[block]
    {\relax\ifthenelse{\NOT\boolean{@tufte@symmetric}}{\begin{fullwidth}}{}}
    {}
    {0pt}
    {\huge\rmfamily\itshape}
    [\ifthenelse{\NOT\boolean{@tufte@symmetric}}{\end{fullwidth}}{}%
      % The fullwidth list must not overwrite the chapter's paragraph hook.
      % Otherwise the first section can lose its normal space after prose.
      \@endpefalse]
  \renewcommand{\maketitlepage}{%
    \begingroup
    \setlength{\parindent}{0pt}
    {\fontsize{24}{24}\selectfont\textit{\@author}\par}
    \vspace{1.75in}
    \begin{fullwidth}
      \raggedright
      {\fontsize{36}{44}\selectfont\@title\par}
    \end{fullwidth}
    \vspace{0.5in}
    {\fontsize{14}{14}\selectfont\textsf{\smallcaps{\@date}}\par}
    \vfill
    {\fontsize{14}{14}\selectfont\textit{\@publisher}\par}
    \thispagestyle{empty}
    \endgroup
  }
  \titlecontents{part}
    [0pt]
    {\addvspace{0.75\baselineskip}\normalfont\normalsize}
    {\allcaps{Part~\thecontentslabel}\quad\allcaps}
    {\allcaps}
    {}
    [\vspace{0.5\baselineskip}]
  % Hide chapter numbers in the contents as well.
  \titlecontents{chapter}
    [4em]
    {\normalfont\normalsize}
    {\textit}
    {\textit}
    {\qquad\thecontentspage}
    [\vspace{0.5\baselineskip}]
}{}
\makeatother

% Enable book and handout links after the Tufte contents customizations.
% The classes' nohyper option defers loading; keep sidenotes unlinked.
\makeatletter
\usepackage[hidelinks,hyperfootnotes=false,bookmarksdepth=0]{hyperref}
% Tufte's pre-expanded author can contain fragile formatting commands.
% Let hyperref strip formatting from the unexpanded author for PDF metadata.
\pdfstringdefDisableCommands{\let\thanklessauthor\@author}
% One caption convention for the book and homework handouts. Load after
% hyperref, and explicitly retain Tufte's margin typography and alignment.
\usepackage{caption}
\DeclareCaptionFont{tufte}{\@tufte@caption@font}
\DeclareCaptionJustification{tufte}{\@tufte@caption@justification}
\captionsetup{
  font=tufte,labelfont=bf,labelsep=period,
  justification=tufte,singlelinecheck=false,parindent=0pt,
  % Tufte builds captions in separate margin boxes; link to the caption.
  hypcap=false
}
\makeatother

% Arrange margin content as a page-sized column, keeping captions with figures.
% https://ctan.org/pkg/marginfix
\usepackage{marginfix}
\usepackage{zref-savepos}
\setlength{\marginparpush}{12pt}
\setlength{\marginheightadjustment}{-6pt}
\makeatletter
\@ifclassloaded{tufte-book}{%
  % Decide a break before starting the heading and its margin insertion.
  \AtBeginDocument{%
    \BeforeBeginEnvironment{exercise}{\Needspace{3\baselineskip}}%
  }
  % Solution links need the final heading baseline, rather than marginfix's
  % page-height estimate, which omits later vertical stretching and floats.
  \let\analysis@marginlabel\@empty
  \newcommand{\analysissolutionnote}[2]{%
    \begingroup
      \edef\analysis@marginlabel{analysis-solution-#1}%
      \leavevmode\zsavepos{\analysis@marginlabel}%
      \marginnote{#2}%
    \endgroup
  }
  % Associate the saved position with both allocated note boxes. Other notes
  % clear the association when LaTeX reuses those boxes.
  \pretocmd{\@xympar}{%
    \expandafter\xdef\csname analysis@marginlabel@\number\@marbox\endcsname
      {\analysis@marginlabel}%
    \expandafter\xdef\csname analysis@marginlabel@\number\@currbox\endcsname
      {\analysis@marginlabel}%
  }{}{\PackageError{analysis-notes}{Could not track solution margin boxes}
    {Check the LaTeX margin paragraph definition.}}
  \let\analysis@margin@getypos\MFX@getypos
  % Record the actual top of each margin column, including pages with top floats.
  \patchcmd{\MFX@attachmargin}{\if\MFX@leftmargin}{%
    \zsavepos{analysis-margin-top-\thepage}%
    \if\MFX@leftmargin
  }{}{\PackageError{analysis-notes}{Could not locate the margin column}
    {Check the marginfix margin attachment definition.}}
  \renewcommand{\MFX@getypos}{%
    \analysis@margin@getypos
    \ifcsname analysis@marginlabel@\number\@marbox\endcsname
      \edef\analysis@label{%
        \csname analysis@marginlabel@\number\@marbox\endcsname}%
      \ifx\analysis@label\@empty\else
        \ifnum\zposy{\analysis@label}>0\relax
          \ifnum\zposy{analysis-margin-top-\thepage}>0\relax
            \dimen@\dimexpr\zposy{analysis-margin-top-\thepage}sp
              -\zposy{\analysis@label}sp-\ht\@marbox\relax
            \edef\mfx@ypos{\the\dimen@}%
          \fi
        \fi
      \fi
    \fi
  }
}{}
% Treat offsets as preferred positions, not whitespace inside the measured box.
% Negative internal space can hide occupied height from the margin scheduler.
\renewcommand{\marginnote}[2][0pt]{%
  \ifthenelse{\boolean{@tufte@loadnatbib}}{%
    \let\cite\@tufte@infootnote@cite
  }{}%
  \gdef\@tufte@citations{}%
  \mparshift{#1}%
  \marginpar{%
    \@tufte@marginnote@font\@tufte@marginnote@justification
    \@tufte@margin@par\noindent #2%
  }%
  \@tufte@print@citations
  \ifthenelse{\boolean{@tufte@loadnatbib}}{%
    \let\cite\@tufte@normal@cite
  }{}%
}
\renewenvironment{@tufte@margin@float}[2][-1.2ex]
  {\FloatBarrier
    \def\analysis@margin@offset{#1}%
    \begin{lrbox}{\@tufte@margin@floatbox}%
      \begin{minipage}[t]{\marginparwidth}%
        \@tufte@caption@font
        \def\@captype{#2}%
        \@tufte@caption@justification
        \@tufte@margin@par\noindent
  }
  {\end{minipage}%
    \end{lrbox}%
    \ifdim\dimexpr\ht\@tufte@margin@floatbox+\dp\@tufte@margin@floatbox\relax
        <\dimexpr\textheight+\marginheightadjustment\relax
      \mparshift{\analysis@margin@offset}%
      \marginpar{\usebox{\@tufte@margin@floatbox}}%
    \else
      \PackageError{analysis-notes}{Margin float is taller than the usable margin}
        {Use a main-text figure or table instead; keep labels readable.}%
    \fi
  }
% Keep Tufte's standard height handling for ordinary and full-width floats.
% Their margin captions must not reserve vertical space in the main text.
\makeatother

% Shared notation and environments. Keep topic-specific macros local.
\newcommand{\N}{\mathbb{N}}
\newcommand{\Z}{\mathbb{Z}}
\newcommand{\Q}{\mathbb{Q}}
\newcommand{\R}{\mathbb{R}}
\newcommand{\C}{\mathbb{C}}

% Draft additions awaiting the author's review. Keep ordinary result styles.
% Remove the wrapper after approval to restore the surrounding text colour.
\definecolor{pendingreviewcolor}{HTML}{486581}
\newenvironment{pendingreview}{%
  \par\color{pendingreviewcolor}%
}{\par}
\newcommand{\pendingreviewtext}[1]{\textcolor{pendingreviewcolor}{#1}}

% One convention for named, linked cross-references in every document.
\renewcommand*{\figureautorefname}{Figure}
\renewcommand*{\tableautorefname}{Table}
\renewcommand*{\partautorefname}{Part}
\renewcommand*{\chapterautorefname}{Chapter}
\renewcommand*{\appendixautorefname}{Appendix}
\renewcommand*{\sectionautorefname}{\S}
\renewcommand*{\subsectionautorefname}{\S}
\renewcommand*{\subsubsectionautorefname}{\S}
\renewcommand*{\theoremautorefname}{Theorem}
% Hyperref appends ~, the reference number, and \null to the name.
% Consume that number here to retain the usual equation parentheses.
\def\equationautorefname~#1\null{Equation~(#1)\null}

% Keep submission-specific layout hints out of shared book exercises.
\makeatletter
\newcommand{\handoutonly}[1]{\@ifclassloaded{tufte-handout}{#1}{}}
% The handouts print solutions immediately; only the book needs these links.
\newcommand{\solutionlink}[1]{%
  \@ifclassloaded{tufte-book}{%
    \analysissolutionnote{#1}{%
      \hyperref[sol:#1]{Solution on page~\pageref*{sol:#1}.}}%
  }{}%
}
\makeatother

% Select the original exercise number for the next appendix restatement.
\newcommand{\restatedexercisekey}{}
\newcommand{\solutionheading}[1]{%
  \renewcommand{\restatedexercisekey}{#1}%
}

% Reuse the exercise's theorem styling and original number. Register only
% the solution anchor, not a second exercise label or solution margin link.
\newenvironment{restatedproblem}{%
  \let\restatedproblemlabel\label
  \renewenvironment{exercise}{%
    \begin{restatedexercise}%
    \phantomsection
    \restatedproblemlabel{sol:\restatedexercisekey}%
  }{\end{restatedexercise}}%
  \renewcommand{\label}[1]{}%
  \renewcommand{\solutionlink}[1]{}%
}{\par}

% Publication month and year used on the copyright page.
\newcommand{\monthyear}{%
  \ifcase\month\or January\or February\or March\or April\or May\or June\or
  July\or August\or September\or October\or November\or December\fi
  \space\number\year
}

% One shared counter makes results unambiguous within each section.
% Alias counters let autoref identify the kind of result, not just its number.
\newcommand{\newsharedtheorem}[2]{%
  \newaliascnt{#1}{theorem}%
  \newtheorem{#1}[#1]{#2}%
  \aliascntresetthe{#1}%
  % Follow later changes to the theorem format, including handout numbering.
  \expandafter\renewcommand\csname the#1\endcsname{\thetheorem}%
  \expandafter\providecommand\csname #1autorefname\endcsname{#2}%
}
\theoremstyle{plain}
\newtheorem{theorem}{Theorem}[section]
\newsharedtheorem{lemma}{Lemma}
\newsharedtheorem{proposition}{Proposition}
\newsharedtheorem{corollary}{Corollary}
\newtheorem*{conjecture}{Conjecture}
\theoremstyle{definition}
\newsharedtheorem{definition}{Definition}
\newsharedtheorem{property}{Property}
\newsharedtheorem{example}{Example}
\newsharedtheorem{exercise}{Exercise}
\newtheorem*{restatedexercise}{\autoref{ex:\restatedexercisekey}}
\theoremstyle{remark}
\newsharedtheorem{remark}{Remark}

% Catalogue numbered theorem statements using their existing counters and
% hyperlink anchors. Match the book's figure/table list formatting.
\makeatletter
\@ifclassloaded{tufte-book}{%
  \newcommand{\listoftheorems}{%
    \begingroup
    % Begin directly after the tables, without a blank facing page.
    \let\cleardoublepage\clearpage
    \chapter*{List of Theorems}%
    \@starttoc{loe}%
    \endgroup
  }
  \let\l@theorem\@tufte@lof@line
  \providecommand*{\toclevel@theorem}{2}
}{}
% Scan this patch directly so etoolbox can handle its parameter characters.
\ifdefined\listoftheorems
  \pretocmd{\@begintheorem}{%
    \ifstrequal{#1}{Theorem}{%
      \ifblank{#2}{}{%
        \addcontentsline{loe}{theorem}{\protect\numberline{#2}#3}%
      }%
    }{}%
  }{}{%
    \PackageError{analysis-notes}{Could not register theorem list entries}
      {Check the amsthm definition of \string\@begintheorem.}%
  }
\fi
\makeatother

% BEGIN TUFTE PLACEHOLDER HELPERS
% Adapted from the supplied Tufte-LaTeX sample (Apache-2.0).
% Remove these documentation helpers once the specimens are replaced.
\newcommand{\hangp}[1]{\makebox[0pt][r]{(}#1\makebox[0pt][l]{)}}

%%
% Prints an asterisk that takes up no horizontal space.
% Useful in tabular environments.

%%
% Prints a trailing space in a smart way.


%%
% Some shortcuts for Tufte's book titles.  The lowercase commands will
% produce the initials of the book title in italics.  The all-caps commands
% will print out the full title of the book in italics.
\newcommand{\vdqi}{\textit{VDQI}\xspace}
\newcommand{\ei}{\textit{EI}\xspace}
\newcommand{\ve}{\textit{VE}\xspace}
\newcommand{\be}{\textit{BE}\xspace}
\newcommand{\VDQI}{\textit{The Visual Display of Quantitative Information}\xspace}
\newcommand{\EI}{\textit{Envisioning Information}\xspace}
\newcommand{\VE}{\textit{Visual Explanations}\xspace}
\newcommand{\BE}{\textit{Beautiful Evidence}\xspace}

\newcommand{\TL}{Tufte-\LaTeX\xspace}

% Typesets the font size, leading, and measure in the form of 10/12x26 pc.
\newcommand{\measure}[3]{#1/#2$\times$\unit[#3]{pc}}

% Macros for typesetting the documentation
\newcommand{\hlred}[1]{\textcolor{Maroon}{#1}}% prints in red
\newcommand{\hangleft}[1]{\makebox[0pt][r]{#1}}
\newcommand{\hairsp}{\hspace{1pt}}% hair space
\newcommand{\hquad}{\hskip0.5em\relax}% half quad space
\newcommand{\ie}{\textit{i.\hairsp{}e.}\xspace}
\newcommand{\eg}{\textit{e.\hairsp{}g.}\xspace}
\newcommand{\na}{\quad--}% used in tables for N/A cells
\providecommand{\XeLaTeX}{X\lower.5ex\hbox{\kern-0.15em\reflectbox{E}}\kern-0.1em\LaTeX}
\newcommand{\tXeLaTeX}{\XeLaTeX\index{XeLaTeX@\protect\XeLaTeX}}
% \index{\texttt{\textbackslash xyz}@\hangleft{\texttt{\textbackslash}}\texttt{xyz}}
\newcommand{\tuftebs}{\symbol{'134}}% a backslash in tt type in OT1/T1
\newcommand{\doccmdnoindex}[2][]{\texttt{\tuftebs#2}}% command name -- adds backslash automatically (and doesn't add cmd to the index)
\newcommand{\doccmddef}[2][]{%
  \hlred{\texttt{\tuftebs#2}}\label{template:cmd:#2}%
  \ifthenelse{\isempty{#1}}%
    {% add the command to the index
      \index{#2 command@\protect\hangleft{\texttt{\tuftebs}}\texttt{#2}}% command name
    }%
    {% add the command and package to the index
      \index{#2 command@\protect\hangleft{\texttt{\tuftebs}}\texttt{#2} (\texttt{#1} package)}% command name
      \index{#1 package@\texttt{#1} package}\index{packages!#1@\texttt{#1}}% package name
    }%
}% command name -- adds backslash automatically
\newcommand{\doccmd}[2][]{%
  \texttt{\tuftebs#2}%
  \ifthenelse{\isempty{#1}}%
    {% add the command to the index
      \index{#2 command@\protect\hangleft{\texttt{\tuftebs}}\texttt{#2}}% command name
    }%
    {% add the command and package to the index
      \index{#2 command@\protect\hangleft{\texttt{\tuftebs}}\texttt{#2} (\texttt{#1} package)}% command name
      \index{#1 package@\texttt{#1} package}\index{packages!#1@\texttt{#1}}% package name
    }%
}% command name -- adds backslash automatically
\newcommand{\docopt}[1]{\ensuremath{\langle}\textrm{\textit{#1}}\ensuremath{\rangle}}% optional command argument
\newcommand{\docarg}[1]{\textrm{\textit{#1}}}% (required) command argument
\newenvironment{docspec}{\begin{quotation}\small\ttfamily\parskip0pt\parindent0pt\ignorespaces}{\end{quotation}}% command specification environment
\newcommand{\docenv}[1]{\texttt{#1}\index{#1 environment@\texttt{#1} environment}\index{environments!#1@\texttt{#1}}}% environment name
\newcommand{\docenvdef}[1]{\hlred{\texttt{#1}}\label{template:env:#1}\index{#1 environment@\texttt{#1} environment}\index{environments!#1@\texttt{#1}}}% environment name
\newcommand{\docpkg}[1]{\texttt{#1}\index{#1 package@\texttt{#1} package}\index{packages!#1@\texttt{#1}}}% package name
\newcommand{\doccls}[1]{\texttt{#1}}% document class name
\newcommand{\docclsopt}[1]{\texttt{#1}\index{#1 class option@\texttt{#1} class option}\index{class options!#1@\texttt{#1}}}% document class option name
\newcommand{\docclsoptdef}[1]{\hlred{\texttt{#1}}\label{template:clsopt:#1}\index{#1 class option@\texttt{#1} class option}\index{class options!#1@\texttt{#1}}}% document class option name defined
\newcommand{\docmsg}[2]{\bigskip\begin{fullwidth}\noindent\ttfamily#1\end{fullwidth}\medskip\par\noindent#2}
\newcommand{\docfilehook}[2]{\texttt{#1}\index{file hooks!#2}\index{#1@\texttt{#1}}}
\newcommand{\doccounter}[1]{\texttt{#1}\index{#1 counter@\texttt{#1} counter}}
% END TUFTE PLACEHOLDER HELPERS


% Preserve original question numbers independently of book numbering.
\renewcommand{\thetheorem}{\arabic{theorem}}

\title{Homework 05}
\author[Lucas Barbosa, lc6424]{Lucas Barbosa, \texttt{lc6424}}
\date{Due Friday, October 16, 2026}

\begin{document}
\maketitle

% !TEX root = ../../hw05.tex
% Source: MA-GY-6213-hw5.pdf, page 1, question 1.
\begin{exercise}
  \label{ex:hw05:interval-fixed-point}
  \solutionlink{hw05:interval-fixed-point}
  Let $F(x)$ be a continuous function which maps the interval $[a,b]$
  into itself. Prove that $F$ has a fixed point in $[a,b]$. That is,
  there is some $c\in[a,b]$ such that $F(c)=c$. Use your result to
  prove that for any $n\in\N$ there is some $x\in(0,1)$ such that
  \[
    x=\frac{1}{1+x^n}.
  \]
  Find $x$ for $n=1$.
\end{exercise}

% !TEX root = ../../hw05.tex
\begin{proof}[Solution]
  A fixed point is a point where $F(x)$ and $x$ have the same value. To
  find one, consider their difference,
  \[
    G(x)=F(x)-x.
  \]
  This function is continuous on $[a,b]$. Since $F$ maps $[a,b]$ into
  itself, $F(a)\ge a$ and $F(b)\le b$. Hence
  \[
    G(a)\ge0,\qquad G(b)\le0.
  \]
  If either value is zero, the corresponding endpoint is a fixed point.
  Otherwise $G(a)>0>G(b)$, so the intermediate value theorem gives a
  $c\in(a,b)$ with $G(c)=0$. Thus $F(c)=c$. This also covers a
  one-point interval, since then $F(a)=a$.

  Now fix $n\in\N$ and define
  \[
    F_n(x)=\frac{1}{1+x^n},\qquad x\in[0,1].
  \]
  The denominator is positive on this interval, so $F_n$ is continuous.
  Also $0\le x^n\le1$, which gives
  \[
    \frac12\le F_n(x)\le1.
  \]
  Therefore $F_n$ maps $[0,1]$ into itself. By the result just proved,
  it has a fixed point $c\in[0,1]$. Neither endpoint is fixed, since
  $F_n(0)=1$ and $F_n(1)=1/2$. Consequently $c\in(0,1)$ and
  \[
    c=\frac{1}{1+c^n}.
  \]
  \autoref{fig:hw05:interval-fixed-point} shows the intersection when
  $n=1$.

  \begin{marginfigure}
    \centering
    % !TEX root = ../../problems/hw05.tex
% Shared physical dimensions for interval and number-line diagrams.
\tikzset{
  interval diagram/.style={
    x=1cm,y=6mm,font=\footnotesize,
    color=black,line width=0.5pt,line cap=round,line join=round,
    >={Stealth[length=4pt,width=3pt]}
  },
  interval endpoint/.style={
    circle,draw=black,line width=0.5pt,
    minimum size=4pt,inner sep=0pt,outer sep=0pt
  },
  interval open/.style={interval endpoint,fill=white},
  interval closed/.style={interval endpoint,fill=black},
  interval label/.style={inner sep=1pt}
}

\begin{tikzpicture}[interval diagram,x=28mm,y=28mm]
  \draw[->] (-0.05,0) -- (1.13,0) node[right=4pt] {$x$};
  \draw[->] (0,-0.05) -- (0,1.14) node[above=4pt] {$y$};
  \node[below left=5pt] at (0,0) {$0$};
  \draw[gray] (1,-0.025) -- (1,0.025);
  \draw[gray] (-0.025,1) -- (0.025,1);
  \node[below=6pt] at (1,0) {$1$};
  \node[left=5pt] at (0,1) {$1$};

  \draw[gray,densely dashed] (0,0) -- (1.04,1.04);
  \node at (0.68,1.09) {$y=x$};
  \draw[domain=0:1,samples=80] plot (\x,{1/(1+\x)});
  \node at (0.92,0.76) {$F_1$};
  \fill (0,1) circle[radius=1.6pt];
  \fill (1,0.5) circle[radius=1.6pt];

  \draw[gray,densely dotted] (0.618034,0) -- (0.618034,0.618034)
    -- (0,0.618034);
  \fill (0.618034,0.618034) circle[radius=1.6pt];
  \node[below=6pt] at (0.618034,0) {$c$};
  \node[left=5pt] at (0,0.618034) {$c$};
\end{tikzpicture}

    \caption[A fixed point as a graph intersection.]{%
      For $n=1$, the graph $y=1/(1+x)$ crosses $y=x$
      inside the unit square. At the crossing, the input equals its output.}
    \label{fig:hw05:interval-fixed-point}
  \end{marginfigure}

  When $n=1$, multiplying by $1+c>0$ gives
  \begin{align*}
    c=\frac{1}{1+c}
      &\iff c(1+c)=1\\
      &\iff c^2+c-1=0\\
      &\iff c=\frac{-1+\sqrt5}{2}
        \quad\text{or}\quad c=\frac{-1-\sqrt5}{2}.
  \end{align*}
  Only the first value belongs to $(0,1)$. Hence the required value is
  $c=(\sqrt5-1)/2$.
\end{proof}


% !TEX root = ../../hw05.tex
% Source: MA-GY-6213-hw5.pdf, page 1, question 2.
\begin{exercise}
  \label{ex:hw05:uniform-continuity-boundedness}
  \solutionlink{hw05:uniform-continuity-boundedness}
  Let $f(x)$ be a continuous function over an interval $I$, with finite
  endpoints $a,b$, which may or may not be closed. Prove that if $f(x)$
  is uniformly continuous over $I$, then $f$ is bounded. With this
  result, show that the function
  \[
    f(x)=\frac{1}{1+x},\qquad x\in(-1,1)
  \]
  is continuous but not uniformly continuous.
\end{exercise}

% !TEX root = ../../hw05.tex
\begin{proof}[Solution]

\textbf{Claim.} Let $I\subseteq \mathbb{R}$ be an interval with finite endpoints $a$, $b$. If
$f: I \rightarrow \mathbb{R}$ is uniformly continuous then $f$ is bounded on $I$.

Suppose $f$ is uniformly continuous on $I$. By definition, for every
$\varepsilon > 0$ there is a $\delta > 0$ such that
$$
\forall x,y \in I,\quad |x - y| < \delta \implies |f(x) - f(y)| < \varepsilon.
$$
Choose $\varepsilon = 1$. Then there exists a $\delta > 0$ such that
$$
|x - y| < \delta \implies |f(x) - f(y)| < 1.
$$
Our goal is to show that there exists a finite constant $M > 0$ such that
$$
|f(x)| \le M,\quad \forall x \in I.
$$

If $I$ is empty, the claim is immediate. Otherwise, fix $x_0\in I$ once
and for all.

Since $a$ and $b$ are finite endpoints, we consider
an arbitrary point $x\in I$.
We'll take $N$ steps between $x_0$ and $x$ such that $|x_k - x_{k- 1}| < \delta$
and $|f(x_k) - f(x_{k-1})| < \varepsilon$.

% a figure showing this interval
\autoref{fig:hw05:uniform-continuity-chain} illustrates
the finite chain of small steps.

\begin{marginfigure}
  \centering
  % !TEX root = ../../problems/hw05.tex
% Schematic finite chain; the open-interval case is shown.
% Shared physical dimensions for interval and number-line diagrams.
\tikzset{
  interval diagram/.style={
    x=1cm,y=6mm,font=\footnotesize,
    color=black,line width=0.5pt,line cap=round,line join=round,
    >={Stealth[length=4pt,width=3pt]}
  },
  interval endpoint/.style={
    circle,draw=black,line width=0.5pt,
    minimum size=4pt,inner sep=0pt,outer sep=0pt
  },
  interval open/.style={interval endpoint,fill=white},
  interval closed/.style={interval endpoint,fill=black},
  interval label/.style={inner sep=1pt}
}

\begin{tikzpicture}[interval diagram,x=7mm,y=1mm]
  \draw[gray] (0,0) -- (6,0);
  \node[interval open] at (0,0) {};
  \node[interval open] at (6,0) {};
  \node[above=5pt] at (0,0) {$a$};
  \node[above=5pt] at (6,0) {$b$};

  \draw (0.7,0) -- (5.3,0);
  \node[fill=white,inner sep=2pt] at (3.5,0) {$\cdots$};
  \foreach \position in {0.7,1.7,2.7,4.3,5.3}{
    \fill (\position,0) circle[radius=1.6pt];
  }
  \node[below=5pt] at (0.7,0) {$x_0$};
  \node[below=5pt] at (1.7,0) {$x_1$};
  \node[below=5pt] at (2.7,0) {$x_2$};
  \node[above=5pt] at (4.3,0) {$x_{N-1}$};
  \node[below=5pt] at (5.3,0) {$x_N=x$};

  \draw[<->,gray] (0.7,10) -- (1.7,10)
    node[midway,above=4pt,text=black] {$<\delta$};
  \draw[<->,gray] (4.3,10) -- (5.3,10)
    node[midway,above=4pt,text=black] {$<\delta$};
\end{tikzpicture}

  \caption[Chaining small steps across a bounded interval.]{%
    A schematic chain inside $I$, with the open-interval
    case shown. Every successive gap is less than $\delta$, so each
    corresponding output change is less than $1$. The dots omit
    intermediate steps in the finite chain.}
  \label{fig:hw05:uniform-continuity-chain}
\end{marginfigure}

Over our interval $[a,b]$, if we take $N$ steps,
we can arrange for each to have size at most
$$
\dfrac{b - a}{N}.
$$
For uniform continuity, it is enough to choose $N$ so that
$$
\dfrac{b - a}{N} < \delta \iff \dfrac{b - a}{\delta} < N.
$$
Since $b - a$ and $\delta$ are both fixed, we can select
an integer $N\ge2$ that
works for every pair of points in $I$.

Such an integer exists by the Archimedean property. Set
\[
  x_k=x_0+\frac{k}{N}(x-x_0),\qquad k=0,\ldots,N.
\]
Each $x_k$ lies between $x_0$ and $x$, so it belongs to $I$ because $I$
is an interval. For $k=1,\ldots,N$,
\[
  |x_k-x_{k-1}|=\frac{|x-x_0|}{N}
  \le\frac{b-a}{N}<\delta.
\]
Uniform continuity therefore gives $|f(x_k)-f(x_{k-1})|<1$ at every
step. This construction also works when $x\le x_0$.


For the first two steps, we have
$$
|f(x_2) - f(x_0)| = |f(x_2) - \underbrace{f(x_1) \mathbin{-} f(x_1)}_{0} - f(x_0)|.
$$
Using the triangle inequality $|A| + |B| \ge |A + B|$,
\begin{align*}
|f(x_2) - f(x_1) + f(x_1) - f(x_0)|
&= |[f(x_2) - f(x_1)] + [f(x_1) - f(x_0)]| ,\\
\mathrel{\Longrightarrow}\quad |f(x_2) - f(x_0)| &\le |f(x_2) - f(x_1)| + |f(x_1) - f(x_0)|.
\end{align*}
We can continue adding central terms like $f(x_1)$ that cancel out to generalize
this to
$$
|f(x_N) - f(x_{0})| \le \sum_{k=1}^N |f(x_k) - f(x_{k-1})|.
$$
Since we chose $\varepsilon=1$,
$$
\sum_{k=1}^N |f(x_k) - f(x_{k-1})| \mathrel{<} \varepsilon N \mathrel{=} N.
$$
Hence we've established that
\begin{align*}
|f(x_N) - f(x_0)| \le \sum_{k=1}^N |f(x_k) - f(x_{k-1})| &< N,\\
\mathrel{\Longrightarrow}\quad |f(x_N) - f(x_0)| & < N.
\end{align*}
In this case $x_N = x$, so we add and subtract a $f(x_0)$ term to arrive
at an expression for $|f(x)|$ in terms of $f(x_0)$,
\begin{align*}
|f(x)| &= |f(x) - f(x_0) + f(x_0)|,\\
|f(x)| &= |[f(x) - f(x_0)] + f(x_0)|.
\end{align*}
We can apply the triangle inequality again,
\begin{align*}
|f(x)| &= |[f(x) - f(x_0)] + f(x_0)|,\\
&\mathrel{\le} |f(x) - f(x_0)| + |f(x_0)|.
\end{align*}
Combining our previous result with the new inequality
and setting $M=N+|f(x_0)|>0$,
\begin{align*}
|f(x)| & \le |f(x) - f(x_0)| + |f(x_0)|,\\
|f(x)| &< N + |f(x_0)|,\\
|f(x)| &< M,\qquad \text{for every }x \in I.
\end{align*}
Both $N$ and $x_0$ are fixed, so $M$ does not depend on $x$.
Since $N$ and $|f(x_0)|$ are both finite, $f$ is bounded on $I$. Therefore we have
proved that if $f$ is uniformly continuous over $I$, then it is bounded.

To prove $f(x) = \frac{1}{1 +x}$ is continuous but not uniformly continuous, we can
first prove the case for regular continuity, then
apply the contrapositive of the claim.

% insert diagram of the function itself with special focus on the the limit point
% at x = -1
\autoref{fig:hw05:unbounded-reciprocal} shows the graph
near the excluded endpoint $-1$.

\begin{marginfigure}
  \centering
  % !TEX root = ../../problems/hw05.tex
% The branch is restricted to (-1,1); neither endpoint is included.
% Shared physical dimensions for interval and number-line diagrams.
\tikzset{
  interval diagram/.style={
    x=1cm,y=6mm,font=\footnotesize,
    color=black,line width=0.5pt,line cap=round,line join=round,
    >={Stealth[length=4pt,width=3pt]}
  },
  interval endpoint/.style={
    circle,draw=black,line width=0.5pt,
    minimum size=4pt,inner sep=0pt,outer sep=0pt
  },
  interval open/.style={interval endpoint,fill=white},
  interval closed/.style={interval endpoint,fill=black},
  interval label/.style={inner sep=1pt}
}

\begin{tikzpicture}[interval diagram,x=15mm,y=8mm]
  \draw[gray,densely dashed] (-1,0) -- (-1,5.9);
  \draw[->] (-1.22,0) -- (1.25,0) node[right=4pt] {$x$};
  \draw[->] (0,-0.08) -- (0,6.1) node[above=4pt] {$y$};
  \foreach \height in {1,2,4}{
    \draw[gray] (-0.04,\height) -- (0.04,\height);
    \node[left=5pt] at (0,\height) {$\height$};
  }
  \foreach \position in {-1,1}{
    \draw[gray] (\position,-0.04) -- (\position,0.04);
    \node[below=7pt] at (\position,0) {$\position$};
  }
  \node[below left=5pt] at (0,0) {$0$};

  \draw[<-,domain=-0.82:1,samples=120] plot (\x,{1/(1+\x)});
  \foreach \position/\height in {0/1,-0.5/2,-0.75/4}{
    \fill (\position,\height) circle[radius=1.6pt];
  }
  \node[interval open] at (1,0.5) {};
\end{tikzpicture}

  \caption[An unbounded continuous function on an open interval.]{%
    The graph of $f(x)=1/(1+x)$ on $(-1,1)$ rises
    without bound as $x$ approaches $-1$ from the right. The dashed line
    is the vertical asymptote. The open point at $(1,1/2)$ is also
    excluded from the graph.}
  \label{fig:hw05:unbounded-reciprocal}
\end{marginfigure}

The numerator $1$ and denominator $1+x$ are continuous,
and $1+x>0$ on $(-1,1)$.
The quotient of two continuous functions is continuous wherever the denominator is
not equal to zero, so for $x \in (-1, 1)$ the function $f(x)$ is continuous.

When we analyze the limit as $x \rightarrow -1^+$,
$$
\lim_{x \rightarrow -1^+} \dfrac{1}{1 + x} = +\infty,
$$

This means that there is no finite upper bound for $f$ on $(-1,1)$.
More explicitly, for any $B>0$, choose an integer $n>\max\{B,1\}$.
Then
\[
  x_n=-1+\frac1n\in(-1,1),\qquad
  f(x_n)=\frac{1}{1+x_n}=n>B.
\]
Thus $f$ is unbounded on $(-1,1)$.

By the contrapositive of the above claim, if $f$ is not bounded, then it cannot be
uniformly continuous over $(-1,1)$.

\end{proof}


% !TEX root = ../../hw05.tex
% Source: MA-GY-6213-hw5.pdf, page 1, question 3.
\begin{exercise}
  \label{ex:hw05:odd-polynomial-root}
  \solutionlink{hw05:odd-polynomial-root}
  Show that if the highest degree term of a polynomial is odd powered,
  then it must have a real root.
\end{exercise}

% !TEX root = ../../hw05.tex
\begin{proof}[Solution]
  Let
  \[
    p(x)=a_mx^m+a_{m-1}x^{m-1}+\cdots+a_0,
  \]
  where $m$ is odd and $a_m\ne0$. Dividing by $a_m$ does not change
  the roots, so we work with
  \[
    q(x)=\frac{p(x)}{a_m}
      =x^m+b_{m-1}x^{m-1}+\cdots+b_0.
  \]
  We want to choose two inputs where the leading term is larger in
  absolute value than all the remaining terms together.

  Set $C=\sum_{k=0}^{m-1}|b_k|$. For $|x|\ge1$, every power
  $|x|^k$ with $k\le m-1$ is at most $|x|^{m-1}$. The triangle
  inequality therefore gives
  \[
    \left|\sum_{k=0}^{m-1}b_kx^k\right|
      \le\sum_{k=0}^{m-1}|b_k||x|^k
      \le C|x|^{m-1}.
  \]
  Choose $R>\max\{1,C\}$. Since $m$ is odd, $(-R)^m=-R^m$.
  Applying the bound at $R$ and $-R$, we obtain
  \begin{align*}
    q(R)&\ge R^m-CR^{m-1}
      =R^{m-1}(R-C)>0,\\
    q(-R)&\le -R^m+CR^{m-1}
      =-R^{m-1}(R-C)<0.
  \end{align*}
  The polynomial $q$ is continuous on $[-R,R]$. By the intermediate
  value theorem, there is a $c\in(-R,R)$ such that $q(c)=0$.
  Since $p(c)=a_mq(c)$, we also have $p(c)=0$.
\end{proof}


% !TEX root = ../../hw05.tex
% Source: MA-GY-6213-hw5.pdf, page 1, question 4.
\begin{exercise}
  \label{ex:hw05:continuous-additive-function}
  \solutionlink{hw05:continuous-additive-function}
  Let $f(x)$ be continuous over $\R$ satisfying the additivity condition
  \[
    f(s+t)=f(s)+f(t),\qquad s,t\in\R.
  \]
  Show that there is a constant $a$ such that $f(x)=ax$.
  (Hint: First prove this when $x$ is an integer and identify what $a$
  is. Then prove this when $x$ is rational. Finally pass to limit for
  any real $x$ by using the denseness of $\Q$ in $\R$.)
\end{exercise}

% !TEX root = ../../hw05.tex
\begin{proof}[Solution]
  Following the hint, we first determine $f$ on the integers, then on
  the rationals, and finally on all of $\R$. Set $a=f(1)$.

  Taking $s=t=0$ in the additivity condition gives
  \[
    f(0)=f(0)+f(0),\qquad\text{so }f(0)=0.
  \]
  For any $t\in\R$, we also have
  \[
    0=f(t+(-t))=f(t)+f(-t),
  \]
  and therefore $f(-t)=-f(t)$. For a positive integer $m$, repeated
  addition gives
  \[
    f(m)=f(\underbrace{1+\cdots+1}_{m\text{ terms}})
      =mf(1)=ma.
  \]
  Together with $f(0)=0$ and $f(-m)=-f(m)$, this proves $f(m)=am$
  for every integer $m$.

  Now let $r=m/n$ be rational, where $m\in\Z$ and $n\in\N$.
  Adding $m/n$ to itself $n$ times gives
  \[
    nf\left(\frac mn\right)
      =f\left(\underbrace{\frac mn+\cdots+\frac mn}_{n\text{ terms}}\right)
      =f(m)=am.
  \]
  Dividing by $n$ yields
  \[
    f\left(\frac mn\right)=a\frac mn.
  \]
  Thus $f(r)=ar$ for every $r\in\Q$.

  \begin{pendingreview}
  Let $x\in\R$ be arbitrary.
  By the density of $\Q$ in $\R$, there exists a sequence $(q_n)$
  of rational numbers such that
  \[
    q_n\longrightarrow x.
  \]
  Since $f$ is continuous at $x$, the sequential characterization of
  continuity gives
  \[
    \lim_{n\to\infty}f(q_n)=f(x).
  \]
  But we have already established that $f(q_n)=aq_n$ for every $n$.
  Therefore,
  \[
    \begin{aligned}
      \lim_{n\to\infty}f(q_n)
        &=\lim_{n\to\infty}aq_n\\
        &=a\lim_{n\to\infty}q_n\\
        &=ax.
    \end{aligned}
  \]
  By the uniqueness of limits,
  \[
    f(x)=ax.
  \]
  Since $x\in\R$ was arbitrary, we conclude
  \[
    f(x)=ax,\qquad\forall x\in\R,\quad a=f(1).\qedhere
  \]
  \end{pendingreview}
\end{proof}

\begin{pendingreview}
The central insight is that additivity determines $f$ on all rational
numbers. Because the rationals are dense in the reals, continuity forces
the function to preserve that same linear relationship at every
irrational number.
\end{pendingreview}


% !TEX root = ../../hw05.tex
% Source: MA-GY-6213-hw5.pdf, page 1, question 5.
\begin{exercise}
  \label{ex:hw05:composition-continuity}
  \solutionlink{hw05:composition-continuity}
  Use the definition of a function being continuous at a point, prove
  that if $f(x)$ is continuous at $x_0$, $g(y)$ is continuous at
  $y_0=f(x_0)$, then $h(x)=g(f(x))=(g\circ f)(x)$ is continuous at
  $x_0$ as well.
\end{exercise}

% !TEX root = ../../hw05.tex
\begin{proof}[Solution]
  Let $\varepsilon>0$. Our goal is to make
  $|g(f(x))-g(f(x_0))|<\varepsilon$ by taking $x$ sufficiently close
  to $x_0$. We first find how close $f(x)$ must be to $f(x_0)$.

  Write $y_0=f(x_0)$. Since $g$ is continuous at $y_0$, there is an
  $\eta>0$ such that
  \[
    |y-y_0|<\eta
      \quad\Longrightarrow\quad |g(y)-g(y_0)|<\varepsilon.
  \]
  Since $f$ is continuous at $x_0$, there is a $\delta>0$ such that
  \[
    |x-x_0|<\delta
      \quad\Longrightarrow\quad |f(x)-f(x_0)|<\eta.
  \]
  The choices are made in this order because the tolerance $\eta$ for
  $f$ comes from the continuity of $g$.
  \autoref{fig:hw05:composition-neighbourhoods} illustrates the two steps.

  \begin{marginfigure}
    \centering
    % !TEX root = ../../problems/hw05.tex
% These neighbourhoods are schematic and do not have a common scale.
% Shared physical dimensions for interval and number-line diagrams.
\tikzset{
  interval diagram/.style={
    x=1cm,y=6mm,font=\footnotesize,
    color=black,line width=0.5pt,line cap=round,line join=round,
    >={Stealth[length=4pt,width=3pt]}
  },
  interval endpoint/.style={
    circle,draw=black,line width=0.5pt,
    minimum size=4pt,inner sep=0pt,outer sep=0pt
  },
  interval open/.style={interval endpoint,fill=white},
  interval closed/.style={interval endpoint,fill=black},
  interval label/.style={inner sep=1pt}
}

\begin{tikzpicture}[interval diagram,x=10mm,y=1mm]
  \foreach \height/\centre/\tolerance in {
    0/{x_0}/{\delta},-24/{y_0=f(x_0)}/{\eta},
    -48/{g(y_0)}/{\varepsilon}}{
    \fill[gray!10] (-1.25,{\height-2}) rectangle (1.25,{\height+2});
    \draw[<->,gray] (-1.6,\height) -- (1.6,\height);
    \node[interval open] at (-1.25,\height) {};
    \node[interval open] at (1.25,\height) {};
    \fill (0,\height) circle[radius=1.6pt];
    \node[below=5pt] at (0,\height) {$\centre$};
    \draw[<->] (0,{\height+7}) -- (1.25,{\height+7})
      node[midway,above=4pt] {$\tolerance$};
  }
  % Leave a separate gutter for arrows, away from the centre labels.
  \draw[->] (-1.85,-4) -- (-1.85,-20)
    node[midway,left=4pt] {$f$};
  \draw[->] (-1.85,-28) -- (-1.85,-44)
    node[midway,left=4pt] {$g$};
\end{tikzpicture}

    \caption[Neighbourhoods for a composition.]{%
      Schematic neighbourhoods for a composition.
      First choose $\eta$ to reach the required output tolerance $\varepsilon$.
      Then choose $\delta$ to reach that intermediate tolerance $\eta$.
      The arrows show the order in which inputs pass through the maps.}
    \label{fig:hw05:composition-neighbourhoods}
  \end{marginfigure}

  If $|x-x_0|<\delta$, the second implication places $y=f(x)$ within
  $\eta$ of $y_0$. Applying the first implication with this value of $y$
  gives
  \[
    |h(x)-h(x_0)|
      =|g(f(x))-g(f(x_0))|<\varepsilon.
  \]
  This is the definition of continuity of $h=g\circ f$ at $x_0$.
\end{proof}


% !TEX root = ../../hw05.tex
% Source: MA-GY-6213-hw5.pdf, page 1, question 6.
\begin{exercise}
  \label{ex:hw05:absolute-value-continuity}
  Use the definition of a function being continuous at a point, show
  that if $f(x)$ is continuous at $x_0$, so is the function
  \[
    g(x)=|f(x)|.
  \]
  \solutionlink{hw05:absolute-value-continuity}
\end{exercise}

% !TEX root = ../../hw05.tex
\begin{proof}[Solution]
  Let $\varepsilon>0$. Since $f$ is continuous at $x_0$, there is a
  $\delta>0$ such that
  \[
    |x-x_0|<\delta
      \quad\Longrightarrow\quad |f(x)-f(x_0)|<\varepsilon.
  \]
  We can use the same $\delta$ for $g(x)=|f(x)|$. The reason is the
  reverse triangle inequality,
  \[
    \bigl||u|-|v|\bigr|\le |u-v|,\qquad u,v\in\R.
  \]
  Taking absolute values cannot increase the distance between two
  outputs, as illustrated in \autoref{fig:hw05:absolute-value-fold}.

  \begin{marginfigure}
    \centering
    % !TEX root = ../../problems/hw05.tex
% The two number lines use the same scale; u < 0 < v < |u|.
% Shared physical dimensions for interval and number-line diagrams.
\tikzset{
  interval diagram/.style={
    x=1cm,y=6mm,font=\footnotesize,
    color=black,line width=0.5pt,line cap=round,line join=round,
    >={Stealth[length=4pt,width=3pt]}
  },
  interval endpoint/.style={
    circle,draw=black,line width=0.5pt,
    minimum size=4pt,inner sep=0pt,outer sep=0pt
  },
  interval open/.style={interval endpoint,fill=white},
  interval closed/.style={interval endpoint,fill=black},
  interval label/.style={inner sep=1pt}
}

\begin{tikzpicture}[interval diagram,x=8mm,y=1mm]
  \draw[<->,gray] (-2.4,0) -- (2.4,0);
  \foreach \position/\name in {-2/u,0/0,1/v}{
    \node[interval closed] at (\position,0) {};
    \node[below=5pt] at (\position,0) {$\name$};
  }
  \draw[<->] (-2,8) -- (1,8)
    node[midway,above=4pt] {$|u-v|$};

  \draw[->,gray] (-2,-7)
    .. controls (-2,-17) and (2,-16) .. (2,-25);
  \node[anchor=west] at (-2.35,-21) {fold at $0$};

  \draw[->,gray] (-0.3,-30) -- (2.4,-30);
  \foreach \position/\name in {0/0,1/{|v|},2/{|u|}}{
    \node[interval closed] at (\position,-30) {};
    \node[below=5pt] at (\position,-30) {$\name$};
  }
  \draw[<->] (1,-41) -- (2,-41)
    node[midway,below=4pt] {$\bigl||u|-|v|\bigr|$};
\end{tikzpicture}

    \caption[Absolute value as a fold of the real line.]{%
      Absolute value folds the negative half-line onto the
      positive half-line. In the illustrated case $u<0<v<|u|$, the
      original distance is $|u|+v$, while the folded distance is
      $|u|-v$. Folding brings the two points closer together.}
    \label{fig:hw05:absolute-value-fold}
  \end{marginfigure}

  Applying the inequality with $u=f(x)$ and $v=f(x_0)$ gives
  \[
    |g(x)-g(x_0)|
      =\bigl||f(x)|-|f(x_0)|\bigr|
      \le |f(x)-f(x_0)|<\varepsilon
  \]
  whenever $|x-x_0|<\delta$. Hence $g$ is continuous at $x_0$.
\end{proof}


% !TEX root = ../../hw05.tex
% Source: MA-GY-6213-hw5.pdf, page 2, question 7.
\begin{exercise}
  \label{ex:hw05:square-not-uniformly-continuous}
  Show that the function $f(x)=x^2$ is not uniformly continuous over
  $\R$ and explain what goes ``wrong''.
  \solutionlink{hw05:square-not-uniformly-continuous}
\end{exercise}

% !TEX root = ../../hw05.tex
\begin{proof}[Solution]
  To disprove uniform continuity, we need one fixed output tolerance
  that no single input tolerance can guarantee over all of $\R$.
  Choose $\varepsilon_0=1$.

  Let $\delta>0$ be arbitrary. By the Archimedean property, choose
  $n\in\N$ with $n>1/\delta$, and set
  \[
    x_n=n,\qquad y_n=n+\frac1n.
  \]
  These inputs are less than $\delta$ apart, since
  \[
    |y_n-x_n|=\frac1n<\delta.
  \]
  However, their outputs satisfy
  \begin{align*}
    |f(y_n)-f(x_n)|
      &=\left(n+\frac1n\right)^2-n^2\\
      &=n^2+2+\frac1{n^2}-n^2\\
      &=2+\frac1{n^2}>1=\varepsilon_0.
  \end{align*}
  Thus every $\delta>0$ fails for the same $\varepsilon_0=1$, so $f$
  is not uniformly continuous on $\R$.

  What goes wrong is visible in the factorisation
  \[
    |y^2-x^2|=|y-x||y+x|.
  \]
  Even when $|y-x|$ is small, the factor $|y+x|$ can be arbitrarily
  large. In our construction the inputs get closer together as they
  move to the right, while the output gap stays above $2$.
  \autoref{fig:hw05:square-output-gap} shows one pair from this sequence.

  \begin{marginfigure}
    \centering
    % !TEX root = ../../problems/hw05.tex
% A magnified local view of x^2 at x_n=3 and y_n=3+1/3.
% Axis origins lie outside this view; horizontal and vertical scales differ.
% Shared physical dimensions for interval and number-line diagrams.
\tikzset{
  interval diagram/.style={
    x=1cm,y=6mm,font=\footnotesize,
    color=black,line width=0.5pt,line cap=round,line join=round,
    >={Stealth[length=4pt,width=3pt]}
  },
  interval endpoint/.style={
    circle,draw=black,line width=0.5pt,
    minimum size=4pt,inner sep=0pt,outer sep=0pt
  },
  interval open/.style={interval endpoint,fill=white},
  interval closed/.style={interval endpoint,fill=black},
  interval label/.style={inner sep=1pt}
}

\begin{tikzpicture}[interval diagram,x=60mm,y=9mm]
  \draw[->] (2.9,8.6) -- (3.5,8.6) node[right=4pt] {$x$};
  \draw[->] (2.9,8.6) -- (2.9,12.2) node[above=4pt] {$y$};
  \draw[<->,domain=2.98:3.43,samples=80] plot (\x,{\x*\x});

  \draw[gray,densely dashed] (3,8.6) -- (3,9) -- (3.333333,9);
  \draw[gray,densely dashed] (3.333333,8.6) -- (3.333333,11.111111);
  \draw[gray] (3.333333,9) -- (3.44,9);
  \draw[gray] (3.333333,11.111111) -- (3.44,11.111111);
  \fill (3,9) circle[radius=1.6pt];
  \fill (3.333333,11.111111) circle[radius=1.6pt];
  \node[below=6pt] at (3,8.6) {$x_n$};
  \node[below=6pt] at (3.333333,8.6) {$y_n$};

  \draw[<->] (3,7.3) -- (3.333333,7.3)
    node[midway,below=4pt] {$1/n$};
  \draw[<->] (3.44,9) -- (3.44,11.111111)
    node[midway,right=4pt] {$>2$};
\end{tikzpicture}

    \caption[Small input gaps with large output gaps.]{%
      A local view of $y=x^2$ for the pair $x_n=n$,
      $y_n=n+1/n$ when $n=3$. A horizontal gap of $1/n$ produces a
      vertical gap of $2+1/n^2$. The two axes use different scales.}
    \label{fig:hw05:square-output-gap}
  \end{marginfigure}
\end{proof}


% !TEX root = ../../hw05.tex
% Source: MA-GY-6213-hw5.pdf, page 2, question 8, parts (a)--(b).
\begingroup
\renewcommand{\labelenumi}{(\alph{enumi})}
\begin{exercise}
  \label{ex:hw05:lipschitz-continuity}
  Consider a function $f(x)$ over $\R$. We say that $f(x)$ is
  Lipschitzian if there is a constant $L>0$ such that
  \[
    |f(x)-f(y)|\le L|x-y|,\qquad x,y\in\R.
  \]
  Here $L$ is called the Lipschitz constant of $f$.
  \begin{enumerate}
    \item Prove that a Lipschitz function is uniformly continuous over
      the entire $\R$.
    \item Consider the function
      \[
        f(x)=\frac{1}{1+x^2},\qquad x\in\R.
      \]
      Show that this function is Lipschitzian and estimate its
      Lipschitzian constant.
  \end{enumerate}
  \solutionlink{hw05:lipschitz-continuity}
\end{exercise}
\endgroup

% !TEX root = ../../hw05.tex
\begin{proof}[Solution]
  \textit{(a)} Let $\varepsilon>0$ and choose $\delta=\varepsilon/L$.
  Since $L>0$, we have $\delta>0$. For any $x,y\in\R$,
  \[
    |x-y|<\delta
      \quad\Longrightarrow\quad
      |f(x)-f(y)|\le L|x-y|<L\delta=\varepsilon.
  \]
  The same $\delta$ works for every pair $x,y$, so $f$ is uniformly
  continuous on $\R$.

  \textit{(b)} For $f(x)=1/(1+x^2)$, we first subtract the fractions
  and factor the numerator,
  \begin{align*}
    |f(x)-f(y)|
      &=\left|\frac{1}{1+x^2}-\frac{1}{1+y^2}\right|\\
      &=\frac{|y^2-x^2|}{(1+x^2)(1+y^2)}\\
      &=|x-y|\frac{|x+y|}{(1+x^2)(1+y^2)}.
  \end{align*}
  It remains to bound the last fraction independently of $x$ and $y$.
  For any $t\in\R$, the inequality $(|t|-1)^2\ge0$ gives
  \[
    t^2-2|t|+1\ge0,
    \qquad\text{so}\qquad |t|\le\frac{1+t^2}{2}.
  \]
  Applying this to both $x$ and $y$, we obtain
  \begin{align*}
    \frac{|x+y|}{(1+x^2)(1+y^2)}
      &\le\frac{|x|+|y|}{(1+x^2)(1+y^2)}\\
      &\le\frac{(1+x^2)+(1+y^2)}{2(1+x^2)(1+y^2)}\\
      &=\frac{1}{2(1+y^2)}+\frac{1}{2(1+x^2)}\\
      &\le\frac12+\frac12=1.
  \end{align*}
  Therefore
  \[
    |f(x)-f(y)|\le |x-y|,\qquad x,y\in\R.
  \]
  Thus $f$ is Lipschitz with $L=1$. This is a valid estimate; we do
  not need the smallest possible constant. By part (a), $f$ is also
  uniformly continuous on $\R$.
\end{proof}


% !TEX root = ../../hw05.tex
% Source: MA-GY-6213-hw5.pdf, page 2, question 9, parts (a)--(b).
\begingroup
\renewcommand{\labelenumi}{(\alph{enumi})}
\begin{exercise}
  \label{ex:hw05:pointwise-not-uniform}
  Consider the sequence
  \[
    f_n(x)=\frac{nx}{1+n^2x^2},\qquad
    x\in[0,1],\quad n=1,2,\ldots.
  \]
  \begin{enumerate}
    \item Show that $f_n(x)\to0$ pointwise as $n\to\infty$.
    \item Show that the convergence is not uniform.
  \end{enumerate}
  \solutionlink{hw05:pointwise-not-uniform}
\end{exercise}
\endgroup

% !TEX root = ../../hw05.tex
\begin{proof}[Solution]
  \textit{(a)} First consider $x=0$. For every $n\in\N$,
  \[
    f_n(0)=0,
  \]
  so the limit at this input is zero. Now fix $x\in(0,1]$. Since
  $1+n^2x^2\ge n^2x^2>0$,
  \[
    0\le f_n(x)=\frac{nx}{1+n^2x^2}
      \le\frac{nx}{n^2x^2}=\frac{1}{nx}.
  \]
  With $x>0$ fixed, $1/(nx)\to0$. The squeeze theorem therefore
  gives $f_n(x)\to0$. Together with the case $x=0$, this proves
  pointwise convergence to the zero function on $[0,1]$.

  \textit{(b)} Pointwise convergence keeps $x$ fixed while $n$ grows.
  To test uniform convergence, we may instead look for a different
  troublesome input at each $n$. Set $x_n=1/n\in[0,1]$. Then
  \[
    f_n(x_n)=\frac{n(1/n)}{1+n^2(1/n)^2}=\frac12.
  \]
  Choose $\varepsilon_0=1/4$. For any proposed $N\in\N$, take
  $n\ge N$ and $x=x_n$. We have
  \[
    |f_n(x)-0|=\frac12>\varepsilon_0.
  \]
  Hence there is no $N$ after which all the values of $f_n$ are
  within $\varepsilon_0$ of zero. The convergence is not uniform.

  We can also see exactly how large the largest error is. Since
  $(nx-1)^2\ge0$, we have $2nx\le1+n^2x^2$ and therefore
  \[
    0\le f_n(x)\le\frac12,\qquad x\in[0,1].
  \]
  Equality holds at $x=1/n$, so
  \[
    \sup_{x\in[0,1]}|f_n(x)-0|=\frac12
  \]
  for every $n$. As \autoref{fig:hw05:moving-peaks} shows, the peak
  moves towards zero while its height stays fixed. This does not
  contradict pointwise convergence, since a fixed positive input
  does not follow the moving peak, and $f_n(0)$ is always zero.

  \begin{marginfigure}
    \centering
    % !TEX root = ../../problems/hw05.tex
% Shared physical dimensions for interval and number-line diagrams.
\tikzset{
  interval diagram/.style={
    x=1cm,y=6mm,font=\footnotesize,
    color=black,line width=0.5pt,line cap=round,line join=round,
    >={Stealth[length=4pt,width=3pt]}
  },
  interval endpoint/.style={
    circle,draw=black,line width=0.5pt,
    minimum size=4pt,inner sep=0pt,outer sep=0pt
  },
  interval open/.style={interval endpoint,fill=white},
  interval closed/.style={interval endpoint,fill=black},
  interval label/.style={inner sep=1pt}
}

\begin{tikzpicture}[interval diagram,x=33mm,y=46mm]
  \draw[->] (-0.03,0) -- (1.1,0) node[right=4pt] {$x$};
  \draw[->] (0,-0.02) -- (0,0.66) node[above=4pt] {$y$};
  \node[below left=5pt] at (0,0) {$0$};
  \draw[gray] (1,-0.015) -- (1,0.015);
  \node[below=6pt] at (1,0) {$1$};
  \draw[gray,densely dotted] (0,0.5) -- (0.62,0.5);
  \node[left=5pt] at (0,0.5) {$\frac12$};

  \draw[domain=0:1,samples=140] plot (\x,{2*\x/(1+4*\x*\x)});
  \draw[gray,domain=0:1,samples=140] plot (\x,{4*\x/(1+16*\x*\x)});
  \draw[densely dashed,domain=0:1,samples=160]
    plot (\x,{8*\x/(1+64*\x*\x)});
  \foreach \position in {0.5,0.25,0.125}{
    \fill (\position,0.5) circle[radius=1.6pt];
  }
  \fill (0,0) circle[radius=1.6pt];

  % The legend has its own space below the plot and its axis labels.
  \draw (0.05,-0.26) -- (0.27,-0.26);
  \node[anchor=west] at (0.35,-0.26) {$n=2$};
  \draw[gray] (0.05,-0.42) -- (0.27,-0.42);
  \node[anchor=west] at (0.35,-0.42) {$n=4$};
  \draw[densely dashed] (0.05,-0.58) -- (0.27,-0.58);
  \node[anchor=west] at (0.35,-0.58) {$n=8$};
\end{tikzpicture}

    \caption[Moving peaks prevent uniform convergence.]{%
      The graphs for $n=2,4,8$ have peaks at $(1/n,1/2)$.
      The peaks move towards the origin without losing height.
      Uniform convergence to zero would require the entire graph
      eventually to lie in every horizontal band about the $x$-axis.}
    \label{fig:hw05:moving-peaks}
  \end{marginfigure}
\end{proof}


% !TEX root = ../../hw05.tex
% Source: MA-GY-6213-hw5.pdf, page 2, question 10, parts (a)--(b).
\begingroup
\renewcommand{\labelenumi}{(\alph{enumi})}
\begin{exercise}
  \label{ex:hw05:coercive-minimum}
  Let $f(x)$ be continuous over $\R$ and satisfy the ``coercive''
  condition:
  \[
    f(x)\ge a|x|-b,\qquad x\in\R,
  \]
  where $a>0$ and $b>0$ are two constants.
  \begin{enumerate}
    \item Show that $\inf f$ exists.
    \item Show that $f$ attains $\inf f$ somewhere at $x_0\in\R$.
      That is, there is some $x_0$ such that
      \[
        f(x_0)=\inf f.
      \]
      Such a point $x_0$ is called a global minimum point of $f$ over $\R$.
  \end{enumerate}
  (Hint: Let $\{x_n\}$ be a sequence such that
  \[
    f(x_n)\to\inf f\quad\text{as }n\to\infty.
  \]
  Show that $\{x_n\}$ contains a convergent subsequence.)
  \solutionlink{hw05:coercive-minimum}
\end{exercise}
\endgroup

% !TEX root = ../../hw05.tex
\begin{proof}[Solution]
  \textit{(a)} Since $a|x|\ge0$ for every $x\in\R$, the given bound
  implies
  \[
    f(x)\ge a|x|-b\ge-b.
  \]
  Thus the range $f(\R)$ is bounded below. It is nonempty because it
  contains $f(0)$. By completeness of $\R$, its infimum exists as a
  real number. Write
  \[
    m=\inf f(\R),\qquad -b\le m\le f(0).
  \]

  \textit{(b)} We follow the hint by choosing inputs whose outputs
  approach $m$. By the defining property of the infimum, for every
  $n\in\N$ there is an $x_n\in\R$ such that
  \[
    m\le f(x_n)<m+\frac1n.
  \]
  Indeed, if no such point existed, $m+1/n$ would be a lower bound
  for $f(\R)$ larger than $m$. These inequalities give $f(x_n)\to m$.

  We next show that the inputs $x_n$ stay in a bounded interval.
  Combining the coercive bound with the choice of $x_n$ gives
  \begin{align*}
    a|x_n|-b&\le f(x_n)<m+\frac1n\le m+1,\\
    a|x_n|&<m+b+1,\\
    |x_n|&<\frac{m+b+1}{a}.
  \end{align*}
  Since $m\ge-b$ and $a>0$, the number
  $R=(m+b+1)/a$ is positive and finite. Thus $(x_n)$ is contained
  in $[-R,R]$. The coercive bound prevents inputs with nearly
  minimal outputs from escaping to infinity, as illustrated in
  \autoref{fig:hw05:coercive-minimum}.

  \begin{marginfigure}
    \centering
    % !TEX root = ../../problems/hw05.tex
% An exact example, f(x)=1+x^2 >= |x|-1. Here m=1 and R=3.
% Shared physical dimensions for interval and number-line diagrams.
\tikzset{
  interval diagram/.style={
    x=1cm,y=6mm,font=\footnotesize,
    color=black,line width=0.5pt,line cap=round,line join=round,
    >={Stealth[length=4pt,width=3pt]}
  },
  interval endpoint/.style={
    circle,draw=black,line width=0.5pt,
    minimum size=4pt,inner sep=0pt,outer sep=0pt
  },
  interval open/.style={interval endpoint,fill=white},
  interval closed/.style={interval endpoint,fill=black},
  interval label/.style={inner sep=1pt}
}

\begin{tikzpicture}[interval diagram,x=6mm,y=4.5mm]
  \draw[->] (-3.4,0) -- (3.5,0) node[right=4pt] {$x$};
  \draw[->] (0,-1.5) -- (0,8.1) node[above=4pt] {$y$};
  \node[below left=5pt] at (0,0) {$0$};

  \draw[gray,densely dashed] (-3,0) -- (-3,2);
  \draw[gray,densely dashed] (3,0) -- (3,2);
  \draw[gray,densely dotted] (-3.15,2) -- (3.15,2);
  \node[above=4pt] at (2.25,2) {$m+1$};
  \node[below=6pt] at (-3,0) {$-R$};
  \node[below=6pt] at (3,0) {$R$};

  \draw[<->,gray] (-3.3,2.3) -- (0,-1) -- (3.3,2.3);
  \draw[<->,domain=-2.55:2.55,samples=100]
    plot (\x,{1+\x*\x});
  \foreach \position in {0.85,0.55,0.25}{
    \fill (\position,{1+\position*\position}) circle[radius=1.6pt];
  }
  \node[interval closed] at (0,1) {};
  \node[left=6pt] at (0,1) {$m$};

  \draw (-2.8,-3.6) -- (-1.9,-3.6);
  \node[anchor=west] at (-1.5,-3.6) {$f(x)=1+x^2$};
  \draw[gray] (-2.8,-5.3) -- (-1.9,-5.3);
  \node[anchor=west] at (-1.5,-5.3) {$|x|-1$};
\end{tikzpicture}

    \caption[A coercive bound traps a minimising sequence.]{%
      An example with $f(x)=1+x^2$ and the lower bound
      $|x|-1$. Here $m=1$. If $f(x)<m+1=2$, the lower bound forces
      $|x|<3$. The indicated points approach the minimum inside
      this bounded interval.}
    \label{fig:hw05:coercive-minimum}
  \end{marginfigure}

  By the Bolzano--Weierstrass theorem, there is a subsequence
  $(x_{n_k})$ converging to some $x_0\in\R$. Continuity of $f$ at
  $x_0$ gives
  \[
    f(x_{n_k})\longrightarrow f(x_0).
  \]
  But the same subsequence satisfies
  \[
    m\le f(x_{n_k})<m+\frac1{n_k},
  \]
  so $f(x_{n_k})\to m$ as $n_k\to\infty$. By uniqueness of limits,
  $f(x_0)=m=\inf f(\R)$. Hence $f$ attains its global minimum.
\end{proof}


\end{document}
